\section{A compact potential for partition edits}
\label{sec:pilot}

The partitions change with the posterior counting measure.  We will pay
their edits by decreases of a potential, so we must first bound its increases
when a comparator member moves.  A Lipschitz estimate at each scale would
charge the same movement once per level.  The improvement comes from local
concentration: when almost all nearby mass lies close to the request, a
minimizing integrand has a much smaller positive slope there.  The errors
in this slope estimate sum to one logarithm through ball-mass ratios.

The potential is a minimum of linear functions of the measure.  This
makes filtering inexpensive by concavity.  To control physical movement,
we analyze the minimizing functions themselves: concentration forces an
auxiliary site of nearly unit weight near the request, and relocating
that site bounds its displacement.  These two properties give the required
slope estimate.

We use the true-metric measures from Section~\ref{sec:setup}.  Thus, after
the new request $x$ has been revealed,
\begin{equation}
 \mu^+=\mu^- -\nu+\delta_x,\qquad
 \nu\le\mu^-,\qquad \nu(Y)=1,\qquad
 e_x=\int_Y d(x,y)\,d\nu(y),
 \label{eq:pilot-mover}
\end{equation}
and both counting measures have total mass $k$.  The radii are
$r_j=r_0\tau^{-j}$, $1\le j\le J$, with $\tau\ge2$.  The expectation of
the sum of the quantities $e_x$ is $Q$.

\subsection{The finite-dimensional minimization}

A \emph{template} consists of fixed numbers $0<\sigma\le1$, $R\ge256$,
and $L\ge0$.  Put
\[
 D_*=10R,\qquad D_1=10D_*=100R.
\]
All distances in the definitions of the following scalar functions are
dimensionless:
\begin{align}
 b(v)&=\max\{0,\min\{1,2-v/\sigma\}\},\nonumber\\
 a(v)&=\begin{cases}
       \min\{v^2/\sigma^2,1\},&0\le v\le R,\\
       2-v/R,&R\le v\le2R,\\
       0,&v\ge2R,
       \end{cases}\nonumber\\
 G(v)&=\max\{0,\min\{1,v/D_*-1\}\}.
 \label{eq:pilot-bumps}
\end{align}
In particular, $b$ is one on $[0,\sigma]$ and zero on $[2\sigma,\infty)$;
$a$ is one on $[\sigma,R]$; and $G$ is zero on $[0,D_*]$ and one on
$[2D_*,\infty)$.  Their Lipschitz constants are, respectively,
$1/\sigma$, at most $2/\sigma$, and $1/D_*$.  At radius $r$ abbreviate
$a_s(p)=a(d(s,p)/r)$, and use the same convention for $b_s,G_s$.

Let $\mathcal Z_r$ be the set of vectors $z=(z_s)_{s\in Y}$ satisfying
\begin{equation}
 \begin{split}
 &z_s\ge0,\qquad
   \sum_{s\in Y}(z_s-G_s(p))_+\le1\quad(p\in Y),\\
 &d(s,s')\ge5\sigma r
       \quad\text{whenever }s\ne s'\text{ and }z_sz_{s'}>0.
 \end{split}
 \label{eq:pilot-feasible}
\end{equation}
Here $(u)_+=\max\{u,0\}$.  For $g:Y\to[0,1]$ with
$r\Lip(g)\le L$, define
\begin{align}
 F_z^g(p)
   &=(1+g(p))\left(1-\sum_s z_sb_s(p)\right)
       +256\sum_s z_s(1+g(s))a_s(p),
       \label{eq:pilot-integrand}\\
 \mathcal L_r(\mu,g)
   &=r\min_{z\in\mathcal Z_r}\int_Y F_z^g(p)\,d\mu(p).
       \label{eq:pilot-potential}
\end{align}

The two terms in $F_z^g$ favor a site near concentrated mass.  A site at
$s$ cancels part of the residual cost inside $B(s,\sigma r)$, while its
surcharge grows quadratically with distance from $s$ until that distance
reaches $\sigma r$.  The large surcharge makes a site at intermediate
distance from a concentration expensive.  The factors involving $g$
allow a partition edit to lower the contribution of a site near the
request.  We will measure that decrease using the old minimizing site
vector, with $g$ chosen from the partition structures.

The site vector is an auxiliary object, with no prescribed total weight.
Its capacity constraints bound the total weight of nearby sites by one.
The gradual cutoff $G$ will let us move a site and repair those constraints
by reducing neighboring weights.  The separation condition has a different
role: it makes the positive supports of the residual bumps $b_s$ disjoint,
so that this repair has a cost independent of the number of sites.
The template superscript on $\mathcal L_r$ is suppressed when it is fixed.

\begin{lemma}[Compactness, concavity, and uniform bounds]
\label{lem:pilot-compact}
For every nonnegative measure $\mu$ on the finite space $Y$, the minimum
in \eqref{eq:pilot-potential} exists.  For fixed $g$, the function
$\mu\mapsto\mathcal L_r(\mu,g)$ is concave, and
\begin{equation}
 0\le\mathcal L_r(\mu,g)\le2r\mu(Y).
 \label{eq:pilot-size}
\end{equation}
Every feasible integrand satisfies
\begin{equation}
 0\le F_z^g\le514,\qquad
 r\Lip(F_z^g)\le\Lambda:=L+1026/\sigma.
 \label{eq:pilot-global-bounds}
\end{equation}
Moreover, the positive supports of the functions $b_s$ with $z_s>0$
are pairwise disjoint.
\end{lemma}

\begin{proof}
Testing capacity at $p=s$ gives $0\le z_s\le1$.  For each pair with
$d(s,s')<5\sigma r$, the separation condition is exactly the closed
equation $z_sz_{s'}=0$.  Thus $\mathcal Z_r$ is a closed subset of the
finite cube $[0,1]^Y$, containing zero.  It is compact, and the objective
is continuous.  Although this set need not be convex, a minimum therefore
exists.

At every point $p$, the weights of all centers in $B(p,D_*r)$ sum to at
most one, because $G_s(p)=0$ there.  Every center contributing to $a_s(p)$
or $b_s(p)$ lies in that ball.  It follows that
\[
 0\le1-\sum_s z_sb_s(p)\le1,
 \qquad \sum_s z_sa_s(p)\le1.
\]
These inequalities give the bound $514=2+512$ for the integrand, its
nonnegativity, and the lower bound in \eqref{eq:pilot-size}.  Taking
$z=0$ gives the upper bound for the minimum.

If $d(p,q)\le r$, every center contributing at either endpoint belongs
to $B(p,(2R+1)r)\subset B(p,D_*r)$, so all its weights together sum to
at most one.  The residual term in \eqref{eq:pilot-integrand} consequently
has Lipschitz bound $(L+2/\sigma)/r$, and the surcharge term has bound
$1024/(\sigma r)$.  If $d(p,q)>r$, the bound $514$ on the range already
implies the claimed Lipschitz bound, since $\Lambda\ge1026$.  Finally,
if $b_s(p)$ and $b_{s'}(p)$ are both positive, then
$d(s,s')<4\sigma r$, contrary to separation for two positive weights.

For nonnegative measures $\mu_0,\mu_1$ and $0\le\theta\le1$, every
feasible $z$ satisfies
\[
 r\int F_z^g\,d(\theta\mu_0+(1-\theta)\mu_1)
 \ge\theta\mathcal L_r(\mu_0,g)
       +(1-\theta)\mathcal L_r(\mu_1,g).
\]
Taking the minimum proves concavity.  Convexity of $\mathcal Z_r$ was
not used.
\end{proof}

\subsection{The three local optimality estimates}

Fix, for this template, constants satisfying
\begin{equation}
 0<\gamma\le\min\{\sigma/64,1/(1+L)\},
 \qquad 0<\delta\le1/4112.
 \label{eq:pilot-concentration-constants}
\end{equation}
They may be decreased for the applications below.  For a point $x$ and
a measure $\mu$ with $M:=\mu(B(x,\gamma r))>0$, write
\begin{equation}
 h=\frac{\mu(B(x,D_1r)\setminus B(x,\gamma r))}{M},
 \qquad
 t=\frac1M\int_{B(x,\gamma r)}\frac{d(x,p)}r\,d\mu(p).
 \label{eq:pilot-shell-moment}
\end{equation}
In particular $0\le t\le\gamma$.  The shell ratio $h$ measures mass
outside the inner ball, whereas $t$ measures how tightly the inner mass
is centered at $x$.  Both errors matter: even when $h=0$, the mass in
the inner ball need not lie at $x$.  In this subsection $h$ denotes only
the shell ratio, not a partition tier.

\begin{lemma}[Local structure and positive slope]
\label{lem:pilot-local}
Assume $h\le\delta$, and let $z$ be any minimizer in
\eqref{eq:pilot-potential}.  Then the following assertions hold.
\begin{enumerate}
\item A center $s$ with $\sigma r/4\le d(s,x)\le Rr/2$ has $z_s=0$.
\item There is a unique positive center $s$ with $d(s,x)<\sigma r/4$.
      Its weight $z_s$ and normalized displacement $w=d(s,x)/r$ satisfy
      \begin{equation}
       z_s\ge\tfrac34,\qquad
       1-z_s\le\frac{512\gamma}{\sigma^2}t+514h,
       \qquad w\le(5+L)(h+t).
       \label{eq:pilot-local-structure}
      \end{equation}
\item If $d(x,y)<\sigma r/8$, then
      \begin{equation}
       [F_z^g(x)-F_z^g(y)]_+
          \le C_{\mathrm{loc}}(h+t)\frac{d(x,y)}r,
       \label{eq:pilot-local-slope}
      \end{equation}
      where one admissible constant is
      \begin{equation}
       C_{\mathrm{loc}}
       =(L+1024/\sigma)(514+8/\sigma)
           +\frac{1024(5+L)}{\sigma^2}.
       \label{eq:pilot-local-constant}
      \end{equation}
\end{enumerate}
\end{lemma}

\begin{proof}
We perform the calculations in units of $r$; integrals of $F$ therefore
omit the common factor $r$ in the potential.

\emph{Deleting a middle-distance site.}
For $\sigma/4\le d(s,x)/r\le R/2$ and $p\in B(x,\gamma r)$,
the distance $d(s,p)/r$ is at least $\sigma/4-\gamma$ and less than $R$.
Consequently
\[
 a_s(p)\ge(1/4-\gamma/\sigma)^2\ge(15/64)^2.
\]
The support of $b_s$ is inside $B(x,D_1r)$.  The coefficient of $z_s$
in the integral objective is therefore at least
\[
 256\int a_s\,d\mu-2\int b_s\,d\mu
 \ge\left[\frac{225}{16}-2(1+\delta)\right]M>0.
\]
Decreasing a weight preserves both constraints in
\eqref{eq:pilot-feasible}.  A positive such weight would contradict
minimality, proving the first assertion.

\emph{Inserting a unit site.}
Clear all weights at centers in $B(x,5D_*r)$, and put weight one at $x$.
Separation holds, since every retained center is farther than $5D_*r$.
To check capacity, consider a test point $p$ where the new site's
contribution is positive.  Then $d(p,x)<2D_*r$.  Every old center
contributing positively at that test satisfies $d(s',p)<2D_*r$, because
$z_{s'}\le1$ and $G_{s'}(p)<z_{s'}$.  Such a center lies in
$B(x,4D_*r)$ and has been cleared.  The single new contribution is at
most one.  At other tests the new site contributes zero and weights have
only been removed.  Thus the competitor is feasible.

There is at most one positive old site within $\sigma r/4$ of $x$ by
separation; denote its weight by $Z$, taking $Z=0$ if it does not exist.
On the entire inner ball its $b$ value is one.  Every other positive
site has distance greater than $Rr/2$ from $x$ by the first assertion,
so its $b$ value is zero on the inner ball.  The old residual is thus
$1-Z$ there, and the old integrand is at least $1-Z$.  The new integrand
there is at most $512(d(x,p)/r)^2/\sigma^2$, whose integral is at most
$512M\gamma t/\sigma^2$.

All changed bumps are supported in
$B(x,(5D_*+2R)r)\subset B(x,D_1r)$.  Outside the inner ball the increase
is at most $514$ pointwise, by Lemma~\ref{lem:pilot-compact}.  Minimality
of the old vector gives
\begin{equation}
 1-Z\le512\gamma t/\sigma^2+514h
       \le512\gamma^2/\sigma^2+514\delta\le\tfrac14.
 \label{eq:pilot-deficit-proof}
\end{equation}
Here both last summands are at most $1/8$ by
\eqref{eq:pilot-concentration-constants}.  This proves existence,
uniqueness, the weight bound, and the deficit bound in
\eqref{eq:pilot-local-structure}.

\emph{Relocating that site and repairing capacity.}
Write $s$ for this near site and $z=z_s$.  The case $w=0$ already has
the desired displacement bound.  For $w>0$, move its weight from $s$ to
$x$.  Simultaneously replace every other weight $z_i$ with center in
$B(x,5D_*r)$ by $(z_i-q)_+$, where
\[
 q=w/D_*.
\]
If $s\ne x$, the old weight at $x$ is zero by the just-proved uniqueness.
Every other positive center lies beyond $Rr/2$ from $x$, so the moved
site preserves separation.

For a capacity test $p$, put
\[
 A=(z-G_s(p))_+,\quad A'=(z-G_x(p))_+,
 \quad B=\sum_{i\ne s}(z_i-G_i(p))_+.
\]
When $A'\le A$, capacity can only improve.  Otherwise $A'>0$, and
$A'-A\le q$ by the Lipschitz constant of $G$.  The same distance
calculation as above puts every center contributing to $B$ within
$4D_*r$ of $x$, hence among the reduced centers.  If
$b_i=(z_i-G_i(p))_+$, its contribution decreases by exactly
$\min\{q,b_i\}$.  For any nonnegative finite family $(b_i)$,
\begin{equation}
 \sum_i\min\{q,b_i\}\ge\min\{q,\sum_i b_i\}.
 \label{eq:pilot-capacity-repair}
\end{equation}
Indeed, a summand with $b_i\ge q$ proves the inequality immediately;
if none exists, the left side is $\sum_i b_i$.  If $B\ge q$, the
reductions pay the increase $A'-A$.  If $B<q$, all the other
contributions vanish and the remaining contribution satisfies $A'\le z\le1$.
This verifies every capacity test, with no dependence on the number of
sites.

For $p$ in the inner ball put $u=d(s,p)/r$ and $v=d(x,p)/r$.
Both are less than $\sigma$, so the moved site's old and new $a$
values are $u^2/\sigma^2$ and $v^2/\sigma^2$.  The triangle inequality
gives $u^2\ge w^2-2wv+v^2$: if $v\le w$, square $u\ge w-v$;
if $v>w$, the right side $(w-v)^2$ follows from $u\ge v-w$.
Consequently, using $L\gamma\le1$,
\begin{align*}
 (1+g(s))u^2-(1+g(x))v^2
 &\ge (1+g(s))(w^2-2wv)-Lwv^2\\
 &\ge w^2-(4+L\gamma)wv
 \ge w^2-5wv.
\end{align*}
The near site's $b$ value stays one there; all other $b$ values stay
zero; and reducing other weights can only lower their surcharges.
The integral gain on the inner ball is therefore at least
\begin{equation}
 \frac{256zM}{\sigma^2}(w^2-5wt).
 \label{eq:pilot-relocation-gain}
\end{equation}

For completeness, the pointwise loss outside the inner ball has a
bound proportional to $w$, rather than a constant per changed site.
Moving the near site's center changes its residual contribution by at
most $2zw/\sigma$, and its surcharge by at most
$256z(4/\sigma+L)w$.  Thus its total increase is at most $zK_0w$, where
\[
 K_0=1026/\sigma+256L.
\]
The reduction of other weights can only increase the residual.  Since
their positive $b$ supports are disjoint, the residual increase is at
most $2q$ at every point.  All changes vanish outside
$B(x,(5D_*+2R)r)$.  The shell loss is therefore at most
$Mh(zK_0+2/D_*)w$.  Combining this with
\eqref{eq:pilot-relocation-gain} and minimality, and dividing by the
positive quantity $256zMw/\sigma^2$, gives
\[
 w\le5t+\frac{\sigma^2}{256}
                   \left(K_0+\frac{2}{zD_*}\right)h
   \le5t+(5+L)h\le(5+L)(h+t).
\]
For the middle inequality, use $z\ge1/2$, $\sigma\le1$,
$D_*\ge2560$, and
$1026/256+4/(256D_*)<5$.  This proves the last local structural estimate.

\emph{Positive slope at the request.}
Put $v=d(x,y)/r<\sigma/8$.  The near site's $b$ is one at both points,
and all other positive sites have $b=0$ at both points.  The residual
part of $F_z^g(x)-F_z^g(y)$ is consequently at most $L(1-z)v$.
Writing $u=d(s,y)/r$, the near-site surcharge difference has positive
part at most
\[
 \frac{512z}{\sigma^2}(w^2-u^2)_+
 \le\frac{1024w}{\sigma^2}v,
\]
because $(w^2-u^2)_+\le2w|w-u|$.
Every other center contributing an $a$ value at either endpoint is
within $(2R+\sigma/8)r<D_*r$ of $x$.  Capacity at $x$, including the
near site, bounds the total of these other weights by $1-z$.
Their surcharge difference has absolute value at most
$1024(1-z)v/\sigma$.  Altogether,
\[
 [F_z^g(x)-F_z^g(y)]_+
 \le\left[(L+1024/\sigma)(1-z)
              +1024w/\sigma^2\right]v.
\]
Finally \eqref{eq:pilot-deficit-proof} and $\gamma\le\sigma/64$ give
$1-z\le(514+8/\sigma)(h+t)$, and the displacement bound proves
\eqref{eq:pilot-local-slope} with the stated constant.
\end{proof}

\subsection{Summing the drift over scales}

\begin{lemma}[Pilot drift bound]
\label{lem:pilot-drift}
Fix a template and constants as in
\eqref{eq:pilot-concentration-constants}.  At each radius let $g_j$ be
any function with $0\le g_j\le1$ and $r_j\Lip(g_j)\le L$.
For measures satisfying \eqref{eq:pilot-mover},
\begin{equation}
 \sum_{j=1}^J
 [\mathcal L_{r_j}(\mu^+,g_j)
        -\mathcal L_{r_j}(\mu^-,g_j)]_+
 \le C_{\mathrm{dr}}\ell e_x.
 \label{eq:pilot-drift}
\end{equation}
The constant is independent of $Y$, $k$, $J$, and $\tau\ge2$.
For example, set
\begin{align}
 p_0&=\left\lceil\log_2(D_1/\gamma)\right\rceil,\nonumber\\
 C_{\mathrm{dr}}
   &=\frac{2056}{\gamma}+\frac{8224}{\sigma}
       +\frac{2\Lambda p_0}{\delta\log2}
       +2C_{\mathrm{loc}}
                   \left(\frac{p_0}{\log2}+2\gamma\right).
 \label{eq:pilot-drift-constant}
\end{align}
\end{lemma}

\begin{proof}
Let $z^{(j)}$ minimize the old objective at radius $r_j$ and write
$F_j=F_{z^{(j)}}^{g_j}$.  Evaluating the new minimum at this vector yields
\begin{equation}
 \begin{split}
 &\mathcal L_{r_j}(\mu^+,g_j)-\mathcal L_{r_j}(\mu^-,g_j)\\
 &\hspace{15mm}\le r_j\int_Y(F_j(x)-F_j(y))\,d\nu(y)
 \le r_j\int_Y[F_j(x)-F_j(y)]_+\,d\nu(y).
 \end{split}
 \label{eq:pilot-candidate-drift}
\end{equation}
If $e_x=0$, the probability measure $\nu$ is supported at $x$ and the
two counting measures agree.  Assume henceforth $e_x>0$.

The radii smaller than $2e_x/\gamma$ contribute at most
\begin{equation}
 514\sum_{r_j<2e_x/\gamma}r_j
 \le\frac{2056}{\gamma}e_x,
 \label{eq:pilot-small-scales}
\end{equation}
by $\tau\ge2$ and \eqref{eq:pilot-global-bounds}.
At every remaining radius, domination by $\mu^-$ and the definition of
$e_x$ give
\begin{equation}
 M_j:=\mu^-(B(x,\gamma r_j))
 \ge\nu(B(x,\gamma r_j))
 \ge1-\frac{e_x}{\gamma r_j}\ge\tfrac12.
 \label{eq:pilot-inner-mass}
\end{equation}
This is a consecutive range of coarse radii.  Define $h_j,t_j$ there
as in \eqref{eq:pilot-shell-moment} with $\mu=\mu^-$.

We record both scale sums explicitly.  Set
$A_j=\mu^-(B(x,D_1r_j))=(1+h_j)M_j$.  For $u\ge0$,
\begin{equation}
 \min\{1,u\}\le\frac{\log(1+u)}{\log2}.
 \label{eq:pilot-capped-log}
\end{equation}
For $0\le u\le1$ this follows from the chord between the endpoint
values of the concave function $\log(1+u)$, and for $u\ge1$ it follows
from monotonicity.  Since $\tau^{p_0}\gamma\ge D_1$, the mass
$A_j$ is at most $M_{j-p_0}$ whenever that coarser index is in the range.
Assigning the value $k$ to missing coarser masses gives
\[
 \sum_j\log(A_j/M_j)
 \le\sum_j(\log M_{j-p_0}-\log M_j)
 \le p_0\log(2k).
\]
All interior logarithms cancel; at most $p_0$ boundary differences
remain, each between masses in $[1/2,k]$.  This is also the instance of
Lemma~\ref{lem:scale-logarithms} needed here.  Thus
\begin{equation}
 \sum_j\min\{1,h_j\}\le\frac{p_0}{\log2}\log(2k).
 \label{eq:pilot-shell-sum}
\end{equation}

For the moment sum only, order the same radii increasingly as
$\rho_1<\cdots<\rho_n$, with $\rho_{i+1}=\tau\rho_i$, and set
$m_i=\mu^-(B(x,\gamma\rho_i))$, $m_0=0$, and
$\Delta_i=m_i-m_{i-1}$.  Partitioning the $j$th inner ball into the
successive differences of these balls gives
\[
 t_j\le\gamma\sum_{i\le j}
                 \tau^{i-j}\frac{\Delta_i}{m_j}.
\]
The first difference includes the entire first ball; hence this bound
also accounts for any atom already inside that ball.  Since $m_j\ge m_i$
for $j\ge i$, interchanging two finite sums gives
\begin{align}
 \sum_{j=1}^n t_j
 &\le\gamma\sum_{i=1}^n\frac{\Delta_i}{m_i}
                  \sum_{j=i}^n\tau^{i-j}
 \le2\gamma\sum_{i=1}^n\frac{\Delta_i}{m_i}\nonumber\\
 &\le2\gamma\left(1+\log\frac{m_n}{m_1}\right)
 \le2\gamma(1+\log(2k)).
 \label{eq:pilot-moment-sum}
\end{align}
Indeed the first ratio $\Delta_1/m_1$ is one, and, for $i\ge2$,
$\Delta_i/m_i\le\int_{m_{i-1}}^{m_i}du/u$.  Empty scale ranges require
no estimate.

Call a remaining radius bad if $h_j>\delta$.  Because $\delta<1$,
\eqref{eq:pilot-shell-sum} bounds the number of bad radii by
$p_0\log(2k)/(\delta\log2)$.  The global Lipschitz bound gives at most
$\Lambda e_x$ in \eqref{eq:pilot-candidate-drift} at each one.  At a good
radius apply Lemma~\ref{lem:pilot-local} to points with
$d(x,y)<\sigma r_j/8$.  Their total contribution, summed over good
radii and integrated against $\nu$, is at most
\[
 C_{\mathrm{loc}}e_x
 \left[\frac{p_0}{\log2}\log(2k)
                      +2\gamma(1+\log(2k))\right].
\]
Here $h_j=\min\{1,h_j\}$ on good radii.  For each fixed exceptional
point $y$, with $d(x,y)\ge\sigma r_j/8$, use boundedness instead:
\begin{equation}
 \sum_{r_j\le8d(x,y)/\sigma}514r_j
       \le\frac{8224}{\sigma}d(x,y).
 \label{eq:pilot-far-scales}
\end{equation}
Integrating proves that all exceptional terms together cost at most
$8224e_x/\sigma$.  This bound has no factor counting scales.
Combining the estimates, using $1+\log(2k)\le2\ell$ and $\ell\ge1$,
gives precisely \eqref{eq:pilot-drift} with
\eqref{eq:pilot-drift-constant}.
\end{proof}

\subsection{Filtering before selecting edits}

The following formulation will allow both heavy-center edits and the
summably weighted tier edits to use the same drift budget.  It also fixes
the order of conditioning: the filtering estimate is applied at every
step with the old coefficients, before selecting any edits.

\begin{corollary}[Budget for summably weighted templates]
\label{cor:pilot-edits}
Let templates be indexed by a finite or countable set $I$, with fixed
weights $\omega_\alpha\ge0$.  Each template has fixed admissible
parameters and the drift constant $C_{\mathrm{dr},\alpha}$ of
Lemma~\ref{lem:pilot-drift}.  Assume
\[
 A:=\sum_{\alpha\in I}\omega_\alpha<\infty,
 \qquad
 B:=\sum_{\alpha\in I}\omega_\alpha C_{\mathrm{dr},\alpha}<\infty.
\]
At each step, filter and drift the counting measure while keeping the
previous functions $g_{\alpha,j}$ fixed.  Then edit these functions using
the new history, retaining their template bounds.  Suppose the weighted
total potential drop in the edit stage is at least a nonnegative charge
$c_t$, taking $c_t=0$ on steps without edits.  Then
\begin{equation}
 \E\sum_t c_t
 \le B\ell Q+2kA\sum_{j=1}^Jr_j
 \le B\ell Q+\frac{2kr_0A}{\tau-1}.
 \label{eq:pilot-edit-budget}
\end{equation}
In particular, if $A,B$ are absolute constants, the multiplicative
budget is $O(\ell Q)$ and the additive bound is independent of the input
law and horizon.
\end{corollary}

\begin{proof}
Let $\mu_{t-1}$ be the previous post-step posterior counting measure.
The posterior identity is
\[
 \E[\mu_t^-\mid\mathcal F_{t-1}]=\mu_{t-1}.
\]
For every old, $\mathcal F_{t-1}$-measurable function
$g=g_{t-1,\alpha,j}$, concavity gives
\begin{equation}
 \E[\mathcal L_{r_j}(\mu_t^-,g)\mid\mathcal F_{t-1}]
       \le\mathcal L_{r_j}(\mu_{t-1},g).
 \label{eq:pilot-filtering}
\end{equation}
One can see this directly without any assumption on convexity of site
vectors: choose a minimizing vector at $\mu_{t-1}$, use it as a
competitor for each refined posterior, and take its conditional
expectation.  Its objective is linear in the measure and therefore
returns exactly the old objective.

Write $\Phi_{t-1}$ for the weighted total potential after the preceding
step, $\Phi_t^-$ for its value after filtering at old $g$'s, and
$\widehat\Phi_t$ for its value after drift, still at old $g$'s.  Write
$\Phi_t$ for its value after the edits.  The three established bounds
are
\[
 \E\Phi_t^-\le\E\Phi_{t-1},\qquad
 \widehat\Phi_t-\Phi_t^-\le B\ell e_x,
 \qquad \Phi_t\le\widehat\Phi_t-c_t.
\]
All these potentials are nonnegative and at most
$2kA\sum_jr_j$.  The same statements hold for a countable family:
finite partial sums give the filtering inequality by convergence of
nonnegative bounded sums, and the absolute bound by
$2kr_j\omega_\alpha$ for each potential makes the signed differences
summable.  The drift bound sums because its weighted constants sum to
$B$.

Taking expectations, summing in $t$, and using
$\E\sum_t e_x=Q$ proves
$\E\sum_t c_t\le B\ell Q+\E\Phi_0-\E\Phi_T$.
Dropping the nonnegative terminal value and applying
\eqref{eq:pilot-size} yields \eqref{eq:pilot-edit-budget}.
In particular, \eqref{eq:pilot-filtering} was summed over \emph{all}
steps.  It was never conditioned on the event that a later edit was
selected.
\end{proof}

For use in verifying the edit hypothesis, if $z$ is an old minimizer at
the already updated measure $\mu^+$, then
\begin{equation}
 \mathcal L_r(\mu^+,g^{\rm old})
      -\mathcal L_r(\mu^+,g^{\rm new})
 \ge r\int_Y(F_z^{g^{\rm old}}-F_z^{g^{\rm new}})\,d\mu^+.
 \label{eq:pilot-edit-candidate}
\end{equation}
Thus the next section only needs a lower bound on the gain of the old
candidate.  In the tier application the template is fixed across tiers
and $\omega_h=e^{-h}$ for $h=1,2,4,\ldots$.  Since these tiers form a
subset of the positive integers,
$\sum_h e^{-h}\le\sum_{n\ge1}e^{-n}=1/(e-1)$, so both constants in
Corollary~\ref{cor:pilot-edits} are bounded independently of the number
of tiers.
